% Section 7: Reproducibility

\subsection{One-line reproduction}

\begin{quote}\ttfamily
git clone https://github.com/Flamehaven-Labs/openai-erdos-eq22-reproduction\\
cd openai-erdos-eq22-reproduction\\
pip install -e ".[dev]"\\
python -m erdos\_ant.verify
\end{quote}

\noindent
Expected terminal output:

\begin{quote}\ttfamily
Verdict: PASS\\
Checks: 21/21 passed\\
eq (2.2) exponent excess: 6.2391e-38 (published \textasciitilde 6.24e-38, rel.err 0.0001)\\
(frozen-report mode: tracked evidence not written; pass --write-evidence to refresh)
\end{quote}

\noindent
By default the command does not write to disk; the working tree stays
clean. Release maintainers pass \texttt{--write-evidence} to refresh the
tracked files
\path{reports/verification_result.json} and
\path{reports/verification_report.md}.

The line \texttt{Checks: 21/21 passed} refers to the verifier's
top-level checklist. One of those checks is a subprocess execution of
the $60$-test pytest suite with external pytest plugin autoloading
disabled; running \texttt{pytest -q} directly executes the same unit
tests outside the verifier wrapper.

\subsection{Pinned environment}

\begin{tabularx}{\textwidth}{@{}p{0.30\textwidth}X@{}}
\toprule
Item & Value \\
\midrule
Repository & \url{https://github.com/Flamehaven-Labs/openai-erdos-eq22-reproduction} \\
Version (paper source) & v0.2.5 \\
Python & 3.11 and 3.12 (CI-verified) \\
numpy & $\geq 1.26$ \\
mpmath & $\geq 1.3$ \\
Test count & $60$ unit tests $+$ $21$ verifier checks \\
CI matrix & Ubuntu $\times$ Windows $\times$ Python 3.11 / 3.12 \\
Licence & MIT \\
Per-release DOIs & \texttt{CITATION.cff} and \href{https://github.com/Flamehaven-Labs/openai-erdos-eq22-reproduction/releases}{GitHub Releases} \\
\bottomrule
\end{tabularx}

\subsection{Supplementary local checks}

In addition to the unit suite and verifier, the release process keeps
author-operated local source-level checks for claim framing,
traceability, and repository hygiene. Those checks are recorded as
engineering evidence only. They are not independent review, are not
mathematical peer review, and are not used as a premise for the
correctness of the eq.~(2.2) reproduction. See \S\ref{sec:limits} for
the bound on what local release checks can and cannot attest to.

\subsection{Citation}

A \texttt{CITATION.cff} file in the repository root provides citation
metadata in the Citation File Format v1.2 standard, recognised by
Zenodo, GitHub, and most citation managers. The recommended citation
for a specific tagged release is the Zenodo DOI associated with that
tag.
